\documentclass[11pt,a4paper]{article}

\usepackage[utf8]{inputenc}
\usepackage[T5]{fontenc}
\usepackage[vietnamese]{babel}
\usepackage{lmodern}
\usepackage[margin=2.5cm]{geometry}
\usepackage{amsmath}
\usepackage{booktabs}
\usepackage[hidelinks]{hyperref}

\title{Báo cáo thực nghiệm E1\\
Uncertainty-aware DINOv3 cho Fine-grained Image Retrieval}
\author{}
\date{14 tháng 9 năm 2026}

\begin{document}
\maketitle

\section{Mục tiêu và giả thuyết}

E1 đánh giá liệu uncertainty ở mức ảnh từ Evidential Deep Learning (EDL)
có cải thiện retrieval khi dùng để sắp xếp lại kết quả của frozen DINOv3
ViT-B/16 hay không. So sánh chính được khóa trước khi đánh giá:

\begin{itemize}
    \item \textbf{R0}: CLS token DINOv3, chuẩn hóa L2 và cosine similarity.
    \item \textbf{R1}: giữ nguyên danh sách R0, sau đó sắp xếp lại top-$N$
    candidate theo uncertainty tăng dần.
\end{itemize}

R2 kết hợp cosine với certainty score $c_{\mathrm{cert}}=1-u$ chỉ là
ablation phụ. Evidential embedding dùng vector $\boldsymbol{\alpha}$ không
thuộc E1 và được dành cho E2.

\section{Phương pháp}

Với $K=100$ development classes, Evidential Head nhận frozen CLS feature
$x\in\mathbb{R}^{768}$ và sinh:

\begin{equation}
e_k=\operatorname{softplus}(f_k(x)),\qquad
\alpha_k=e_k+1,\qquad
S=\sum_{k=1}^{K}\alpha_k,\qquad
u=\frac{K}{S}.
\end{equation}

Head được tối ưu bằng expected squared-error Bayes risk và annealed KL
regularization theo Sensoy et al.\footnote{M. Sensoy, L. Kaplan, and
M. Kandemir, \href{https://arxiv.org/abs/1806.01768v3}
{Evidential Deep Learning to Quantify Classification Uncertainty},
arXiv:1806.01768v3, 2018.} R1 thích nghi uncertainty-driven top-$N$
reranking của Đorđević và Kumar\footnote{D. Đorđević and S. Kumar,
\href{https://arxiv.org/abs/2409.01082v2}
{Evidential Transformers for Improved Image Retrieval},
arXiv:2409.01082v2, CC BY 4.0, 2025.}, nhưng thay backbone đã fine-tune
trong bài báo bằng frozen DINOv3. Do đó đây là DINOv3 adaptation, không
phải faithful reproduction.

R2 sử dụng:

\begin{equation}
s'=(1-\beta)\widetilde{s}_{\mathrm{cos}}
   +\beta\widetilde{c}_{\mathrm{cert}},
\end{equation}

trong đó hai thành phần được min--max normalize trong top-$N$.

\section{Thiết lập thực nghiệm}

\begin{itemize}
    \item Dataset: CUB-200-2011; 5,864 development images thuộc classes
    0--99 và 5,924 test images thuộc classes 100--199.
    \item Development split: stratified 80\% fit / 20\% validation theo
    từng lớp; test labels không được dùng để chọn tham số.
    \item Query và gallery là cùng tập; self-match được loại trước top-$K$.
    \item Metrics: Recall@1, @2, @4, @8; ba seed 42, 43 và 44.
    \item R1 validation grid:
    $N\in\{10,20,50,100\}$. Cả ba seed chọn $N=10$.
    \item R2 validation grid:
    $\beta\in\{0.1,0.25,0.5,0.75\}$; cả ba seed chọn $\beta=0.1$.
    Seed 42 chọn $N=50$; seed 43 và 44 chọn $N=10$.
    \item Hardware: 2$\times$ Tesla T4; Python 3.12.13,
    PyTorch 2.10.0+cu128, CUDA 12.8, cuDNN 9.1.0.2.
    \item DINOv3 checkpoint:
    \texttt{facebook/dinov3-vitb16-pretrain-lvd1689m}, resolved revision
    \texttt{5931719e67bbdb9737e363e781fb0c67687896bc}.
\end{itemize}

Checkpoint selection theo validation EDL loss chọn epoch 30 cho seed 42 và
43, epoch 20 cho seed 44. Sau đó head được khởi tạo lại và train trên toàn
bộ development set đúng số epoch đã chọn.

\section{Kết quả}

\subsection{Kết quả retrieval chính}

\begin{table}[h]
\centering
\caption{Recall@K trên test set, mean $\pm$ sample standard deviation
qua ba seed. R0 không có thành phần học nên độ lệch chuẩn bằng 0.}
\begin{tabular}{lrrrr}
\toprule
Phương pháp & R@1 & R@2 & R@4 & R@8 \\
\midrule
R0: cosine
& 87.998 & 92.758 & 95.257 & 97.012 \\
R1: uncertainty reranking
& $80.036\pm0.240$ & $88.364\pm0.169$
& $93.991\pm0.058$ & $96.855\pm0.010$ \\
R2: certainty fusion (phụ)
& $87.970\pm0.039$ & $92.719\pm0.042$
& $95.341\pm0.017$ & $97.007\pm0.035$ \\
\bottomrule
\end{tabular}
\end{table}

So với R0, R1 giảm trung bình 7.962 điểm phần trăm tại R@1, 4.395 tại
R@2, 1.266 tại R@4 và 0.158 tại R@8. R2 gần như giữ nguyên baseline:
$-0.028$, $-0.039$, $+0.084$ và $-0.006$ điểm phần trăm tương ứng.
R2 không được dùng để quyết định giả thuyết chính.

\subsection{Paired bootstrap cho R1 so với R0}

\begin{table}[h]
\centering
\caption{Recall@1 delta và paired bootstrap 95\% CI trên 5,924 query;
2,000 bootstrap samples mỗi seed. Đơn vị là điểm phần trăm.}
\begin{tabular}{lrrr}
\toprule
Seed & Delta & CI thấp & CI cao \\
\midrule
42 & $-7.917$ & $-8.913$ & $-6.938$ \\
43 & $-8.221$ & $-9.166$ & $-7.259$ \\
44 & $-7.748$ & $-8.727$ & $-6.769$ \\
\bottomrule
\end{tabular}
\end{table}

Cả ba khoảng tin cậy đều hoàn toàn dưới 0. Tại R@8, mức giảm nhỏ và CI
cắt 0; tuy nhiên suy giảm tại R@1, R@2 và R@4 là rõ ràng.

\subsection{Validation và thay đổi top-1}

Ngay trên validation, R1 với $N=10$ giảm R@1 từ 83.517 xuống 61.512
(seed 42), từ 83.093 xuống 65.421 (seed 43), và từ 84.027 xuống 62.787
(seed 44). Các $N$ lớn hơn làm kết quả giảm mạnh hơn, vì vậy $N=10$ là
lựa chọn ít bất lợi nhất trong grid chứ không phải một cấu hình vượt R0.

Trên test, R1 sửa đúng lần lượt 226, 220 và 239 query vốn sai, nhưng làm
hỏng 695, 707 và 698 query vốn đúng. Chênh lệch bất lợi này giải thích
trực tiếp mức giảm Recall@1 và nhất quán trên cả ba seed.

\subsection{Chất lượng uncertainty}

\begin{table}[h]
\centering
\caption{Khả năng dùng uncertainty của original top-1 candidate để phát
hiện retrieval error, mean $\pm$ sample standard deviation.}
\begin{tabular}{lrrr}
\toprule
Metric & AUROC & AUPRC & AURC \\
\midrule
Giá trị & $0.560\pm0.023$ & $0.162\pm0.006$ & $0.103\pm0.006$ \\
\bottomrule
\end{tabular}
\end{table}

Tỉ lệ lỗi top-1 của R0 là khoảng 0.120, nên AUPRC 0.162 cao hơn prevalence
nhưng AUROC 0.560 chỉ cho thấy tín hiệu phân biệt yếu. Mean uncertainty
trên toàn test set là 0.976, 0.969 và 0.985 cho ba seed; range tập trung
gần maximum uncertainty. Uncertainty của top-1 candidate sai cao hơn
candidate đúng lần lượt khoảng 0.0082, 0.0069 và 0.0023, nhưng khoảng
tách biệt quá nhỏ để thay thế thứ tự cosine bằng hard uncertainty order.

\section{Kiểm tra artifact}

Feature cache chứa đủ 11,788 ảnh và embedding shape
$[11788,768]$. Mỗi final test artifact chứa 5,924 image IDs duy nhất,
cùng full evidence $\mathbf e\in\mathbb{R}^{5924\times100}$,
$\boldsymbol\alpha\in\mathbb{R}^{5924\times100}$ và
$u\in\mathbb{R}^{5924}$. Kiểm tra trực tiếp xác nhận:

\begin{equation}
\boldsymbol\alpha=\mathbf e+1,\qquad
\mathbf p=\boldsymbol\alpha/S,\qquad
u=100/S,
\end{equation}

và toàn bộ giá trị đều hữu hạn. Các tensor này đủ để E2 tái sử dụng mà
không cần forward lại Evidential Head.

\section{Kết luận}

Theo decision rule đã đăng ký trước, \textbf{giả thuyết chính E1 không
được hỗ trợ}. Paper-faithful R1 không cải thiện R0; ngược lại, nó gây suy
giảm có ý nghĩa tại Recall@1 trên cả ba seed. Kết quả hỗ trợ nhận định rằng
image-level uncertainty của một classifier học trên development classes
không phải là đại diện phù hợp cho pairwise retrieval correctness trên
unseen classes.

R2 cho thấy giữ cosine làm thành phần chi phối ($\beta=0.1$) tránh phần
lớn suy giảm, nhưng không tạo cải thiện Recall@1 ổn định và chỉ được xem
là secondary ablation. Bước tiếp theo hợp lý là giữ nguyên negative result
của E1, tái sử dụng full $(\mathbf e,\boldsymbol\alpha,u)$ trong E2,
và chuyển trọng tâm sang pair-wise confidence ở E2B thay vì tiếp tục
tối ưu hard image-level reranking trên test set.

\section{Giới hạn và provenance}

Run được ghi nhận tại Git commit
\texttt{436375d6f08c4f85b65b1f3039eb4dffe8f7a761}, nhưng metadata đánh
dấu \texttt{git\_dirty=true}. Vì vậy kết quả có thể phân tích nội bộ,
song cần chạy xác nhận từ clean commit trước khi dùng làm kết quả công bố.
Ngoài ra, output hiện không chứa metric của random-reranking control R3
và chưa có qualitative image panels; không suy diễn kết quả cho hai phần
chưa được chạy này.

\end{document}
